%% Cover letter for the revised submission to Acta Astronautica (AA-D-26-02719, revision 1)
%% Compiles as a standalone PDF:  pdflatex cover_letter.tex

\documentclass[11pt,a4paper]{article}

\usepackage[a4paper,top=0.55in,bottom=0.55in,left=0.9in,right=0.9in]{geometry}
\usepackage[T1]{fontenc}
\usepackage{lmodern}
\usepackage{microtype}
\usepackage{parskip}
\usepackage{hyperref}
\hypersetup{colorlinks=true,linkcolor=black,urlcolor=blue,citecolor=black,pdfborder={0 0 0}}

\setlength{\parindent}{0pt}
\setlength{\parskip}{0.32em}

\begin{document}
\pagestyle{empty}

\begin{flushright}\today\end{flushright}

\noindent
The Editor\\
\textit{Acta Astronautica}

\vspace{0.3em}

\noindent\textbf{Subject:} Revised manuscript AA-D-26-02719 ---
\textit{``Eclipse-gated transient energy gain for passive solar sails in near-polar low Earth orbit''}

Dear Editor,

Thank you for the opportunity to revise this manuscript, and please convey my thanks to both reviewers for reports that were exacting and, in the end, very useful. I enclose the revised manuscript, a marked-up copy, and a point-by-point response to all fourteen reviewer comments.

\textbf{What changed, and why.} Reviewer~2's central objection --- that the open-source propagator released with the paper does not reproduce the headline 700~km result --- was correct. In tracing it I audited the separate ``reduced-model'' implementation that had produced the original Sections~6--8 and found three implementation errors in it (an inverted incidence-angle convention in its panel aerodynamics, an SRP coefficient of \(1+C_r\) in place of \(C_r\), and a Sun direction frozen at right ascension~\(0^\circ\)). I have therefore withdrawn every quantitative result of the original Sections~6--8 and regenerated them with the open-source propagator described in Section~9, which is cross-validated against Orekit to \(10^{-10}\)~km and against an analytic drag integral. The same audit found and corrected a smaller defect in the open-source propagator itself --- an atmosphere table whose scale heights were shifted by one layer, which had inflated its 14-day drag decay at 700~km from 70 to 84~km. The manuscript states this openly.

The physics survives the correction. Eclipse timing still flips the sign of the orbit-averaged SRP work, and the gain regime is still narrow, transient and self-limiting, exactly as the paper argued --- but it is realised at 800~km rather than 700~km, with a 14-day gain of \(5.95\)~km rather than 20~km, and 700~km is drag-dominated throughout. The abstract, highlights, figures and tables have been rebuilt around these numbers.

\textbf{The STK validation is now complete rather than partial.} Reviewer~1 asked why STK had applied no SRP force and why its drag decay was an order of magnitude below the open-source value. By reading STK's effective inputs back during propagation I found both causes: a default 50~kg fuel load in the Astrogator initial-state segment that had silently increased the propagated mass from 10 to 60~kg, and a scalar-first quaternion passed to an interface that expects the scalar last, which had mis-oriented the sail plate. With the two corrected, STK reproduces the open-source SRP response and the eclipse-gated gain, and the residual drag difference of \(5\)--\(15\%\) is traced to STK's co-rotating-atmosphere and geodetic-altitude conventions, quantified independently in Orekit. Section~9 now states the exact STK version and configuration, and the response letter documents the diagnosis.

\textbf{Other substantive additions} in response to the reviewers: a quantified treatment of atmosphere co-rotation and geodetic altitude; an atmosphere-model and solar-activity study (Harris--Priester and NRLMSISE-00 at low, moderate and high activity) that gives the regime boundary as a function of the solar cycle; a step-size convergence study motivated by the discontinuous eclipse-gated force; consistent metric definitions and new parameter tables; the corrected SRP force equation; the LTAN convention adopted throughout; and the reference and affiliation points raised by Reviewer~1.

I recognise that this is a larger revision than usual. I believe it is the only honest response to the reviews, and that the corrected paper is stronger for it: every number in it is now produced by one inspectable toolchain, reproduced at its key points by two independent references, and anchored by an analytic check.

The manuscript remains original, is not published or under consideration elsewhere, and I am its sole author.

\vspace{0.25em}
\noindent Yours sincerely,

\vspace{0.7em}
\noindent
Ian Poon\\
Department of Bioengineering, Imperial College London\\
\href{mailto:ip323@ic.ac.uk}{ip323@ic.ac.uk}

\end{document}
